\section{Max-Cut 归约}\label{sec:reduction}

固定有理数 $t\in(0,1)$、树参数 $s,m$、有理数 $0<\mu<1/2$ 与
$r=1-\mu/d$。这些参数将在 \Cref{sec:completion} 中选定。设
$\mathcal G$ 是 \Cref{prop:affine-games} 中的仿射博弈，边律为
$\lambda$，字母表维数为 $n+1,n$。设 $T\ge1$ 为整数。我们使用乘积
标记空间
\[
 \mathcal B=(\F^{n+1})^T,\qquad
 \mathcal A=(\F^n)^T,
\]
以及行空间 $H=\Aff(\mathcal A,\F)$。

\subsection{顶点与边分布}

对每个 $\underline x\in X^T$，为每个词 $w\in\bits^{\mathcal B}$ 建一个
图顶点 $(\underline x,w)$。这是标准的长码字域~\cite{BGS98,Hastad01}；
树构造规定测试如何采样其查询。于是任意割恰是一族函数
\[
 \sigma_{\underline x}:\bits^{\mathcal B}\longrightarrow\bits.
\]
这些函数只接收所得的词；它们不接收生成词的行、门表、切片或采样到的
博弈边。

图的边分布是下述测试。
\begin{enumerate}
\item 从 $\lambda$ 的乘积右端点边际中抽取 $\underline y\in Y^T$。
独立地由给定 $\underline y$ 的乘积条件边律抽取两条元组
$\underline e,\underline e'$。记其左端点为 $\underline x,\underline x'$，
其逐坐标投影为 $\pi_{\underline e},\pi_{\underline e'}:\mathcal B\to\mathcal A$。

\item 与所有博弈数据独立地抽取随机切片 $\mathcal D$。按
\Cref{lem:tree-correlation} 抽取 $F,F'$：在 $\mathcal D$ 上相应表项
为独立的 $t$-相关均匀对，在 $\mathcal D$ 外为相同的均匀对。

\item 抽取均匀的公共行数组 $R\in H^k$。在两个副本中各自独立地以概率
$r$ 保留每行、否则均匀重采样，得到 $R^1,R^2$。

\item 输出以
\begin{equation}\label{eq:test-queries}
 \bigl(\underline x,W(R^1,F)\circ\pi_{\underline e}\bigr),
 \qquad
 \bigl(\underline x',-W(R^2,F')\circ\pi_{\underline e'}\bigr)
\end{equation}
为端点的无序边。
\end{enumerate}
所有边权即为这些结果的概率；相同无序端点的结果被合并聚合。

\begin{lemma}[完备性]\label{lem:completeness}
若 $\Val(\mathcal G)\ge1-a$，则图有值至少为
\begin{equation}\label{eq:completeness}
 (1-2Ta)\frac{1+t r^{2d}}2
 \ge\frac{1+t}{2}-t\mu-2Ta
\end{equation}
的割。第一个下界在 $2Ta\le1$ 时有用，这正是下文使用的范围。
\end{lemma}

\begin{proof}
选取以概率至少 $1-a$ 满足边的博弈标记 $\ell_X,\ell_Y$。用左端标记
元组处的取值为每个词顶点染色：
\[
 \sigma_{\underline x}(w)=
 w\bigl(\ell_X(x_1),\ldots,\ell_X(x_T)\bigr).
\]
采样到的边元组的每个坐标都具有原边律。并界表明：所有 $2T$ 个采样
约束都以至少 $1-2Ta$ 的概率被满足。当所有约束被满足时，两个被查询
的元组投影到同一标记
$l=(\ell_Y(y_1),\ldots,\ell_Y(y_T))\in\mathcal A$。此时割符号之积为
$-W(R^1,F)(l)W(R^2,F')(l)$。词生成数据与博弈边独立，故
\Cref{lem:tree-correlation} 给出条件割概率 $(1+t r^{2d})/2$。在其余
事件上割概率非负。当 $2Ta\le1$ 时即得第一个界；若 $2Ta>1$，该界为
负，自动成立。最后，$r^{2d}\ge1-2\mu$，且前因子带来的损失至多
$2Ta$，得第二个界。
\end{proof}

\subsection{大割产生超额稳定性}

任意割在其被查询的词取反时未必变号。我们把它分解为两个奇偶部，因为
第二个查询前的负号使奇性依赖增加割值、偶性依赖减少割值。对每张割表
定义其偶部与奇部：
\[
 \sigma^{\mathrm{even}}(w)=\frac{\sigma(w)+\sigma(-w)}2,
 \qquad
 \sigma^o(w)=\frac{\sigma(w)-\sigma(-w)}2.
\]
对边元组与右端点元组，令
\begin{align}
 G_{\underline e}(R,F)
 &=\sigma^o_{\underline x(\underline e)}
       (W(R,F)\circ\pi_{\underline e}),\label{eq:edge-function}\\
 G_{\underline y}(R,F)
 &=\E_{\underline e\mid\underline y}G_{\underline e}(R,F),
 \qquad \bar G_{\underline y}=T_rG_{\underline y}.
 \label{eq:average-function}
\end{align}
把用 $\sigma^{\mathrm{even}}$ 所作的类似平均记为
$A_{\underline y}$ 与 $\bar A_{\underline y}=T_rA_{\underline y}$。
这些函数的绝对值以 1 为界。此外，
\[
 G_{\underline y}(R,F)=\phi_{\underline y}(W(R,F)),\qquad
 \phi_{\underline y}(w)=\E_{\underline e\mid\underline y}
 \sigma^o_{\underline x(\underline e)}(w\circ\pi_{\underline e}),
\]
其中 $\phi_{\underline y}$ 是奇的，且不依赖任何切片或行数据。于是由
\Cref{lem:tree-correlation} 之后的观察，对每个固定的
$\underline y,\mathcal D,R,F_{\mathrm{out}}$，$\bar G_{\underline y}$
在切片表项上均值为零。后文的抽取论证（\Cref{prop:extraction}）恰好
适用于这种有界奇性的生成词函数。

\begin{lemma}[稳定性恒等式]\label{lem:cut-stability}
对构造出的图的每个割，
\begin{equation}\label{eq:cut-stability}
 \cut(\sigma)=\frac12+
 \frac12\E\Stab_t^{F_{\mathcal D}}(\bar G_{\underline y})
 -\frac12\E\Stab_t^{F_{\mathcal D}}(\bar A_{\underline y}),
\end{equation}
其中两个期望都取遍 $\underline y,\mathcal D,R,F_{\mathrm{out}}$。
特别地，
\begin{equation}\label{eq:cut-upper-stability}
 \cut(\sigma)\le\frac12+
 \frac12\E\Stab_t^{F_{\mathcal D}}(\bar G_{\underline y}).
\end{equation}
\end{lemma}

\begin{proof}
以 $\underline y,\mathcal D,R,F_{\mathrm{out}}$ 为条件。采样到的边
元组与两次行重采样在两个副本之间独立。对它们取平均后，作为
$F_{\mathcal D}$ 的函数，第一个符号的期望是
$\bar A_{\underline y}+\bar G_{\underline y}$。第二个被查询词前的负
号使第二个符号在 $F'_{\mathcal D}$ 处的期望为
$\bar A_{\underline y}-\bar G_{\underline y}$。两个切片输入是普通的
$t$-相关均匀立方输入。这些函数中的行是公共潜数组 $R$；两次独立的行
重采样已被吸收进 $T_r$。把每个切片表项取反，$\bar A_{\underline y}$
不变而 $\bar G_{\underline y}$ 变号。因此二者的傅里叶支撑的总奇偶性
相反，交叉项消失。期望符号积为
$\Stab_t(\bar A_{\underline y})-\Stab_t(\bar G_{\underline y})$。
一条边以「二分之一减符号积的二分之一」的概率被割，给出
\eqref{eq:cut-stability}。$t>0$ 时稳定性非负给出
\eqref{eq:cut-upper-stability}。
\end{proof}

若某割的值至少为 $(1+B(t))/2+\delta/2$，则其奇平均的期望切片稳定性
至少为 $B(t)+\delta$。接下来三节表明：这一超额稳定性产生一个承载
可观傅里叶质量的标记，且常数与博弈维数及重复次数无关；形式化陈述是
\Cref{prop:extraction}。\Cref{sec:decoding} 再把该质量转换为对原
博弈单个副本的染色。
